\subsection{A simple control certificate remains predictive in a Transformer}
\label{sec:nonlinear-suffix}

Equation~\ref{eq:certificate} is exact for an affine system. A Transformer is
nonlinear. Our first experiment asks whether the same three sources of
error---remaining reported error, starting mismatch, and response mismatch---
remain useful when every edit passes through learned Transformer layers. We
begin with a small model whose target we can verify exactly, following the
controlled ``wind-tunnel'' approach of \cite{agarwal2025bayesian}.

\paragraph{Task contract.}
This experiment instantiates the control mode as follows:

\begin{table}[H]
\centering
\small
\begin{tabular}{@{}L{0.21\textwidth}L{0.71\textwidth}@{}}
\toprule
Dimension & Fixed contract \\
\midrule
Target & Scalar state readout and its error over a 15-step trajectory. \\
Information boundary & The residual after the first Transformer block;
first-order observers use the two individual features, while lifted observers
also use their product. \\
Action space & Repeated residual edits along the declared observer direction. \\
Decision rule and loss & Proportional feedback with $K=0.35$; integrated
squared tracking error is primary, with final error and clipping as diagnostics. \\
Operating regime & Twelve independently trained two-block Transformers for
additive and interactional targets, with edits checked against clean residual
scale. \\
\bottomrule
\end{tabular}
\caption{ObserverBench contract for the closed-loop control experiment.}
\label{tab:control-task-contract}
\end{table}

We train a two-block Transformer on sequences with two binary features, $x_1$
and $x_2$, and target
\[
y=0.35x_1+0.25x_2+\gamma x_1x_2+\epsilon.
\]
The product $x_1x_2$ matters only when both features are present. A first-order
observer reads $x_1$ and $x_2$ separately, so it cannot represent this joint
condition. A lifted observer adds the pairwise feature $x_1x_2$. We edit the
residual after block 0 and rerun block 1, the final layer normalization, and the
target head. The edit therefore passes through the learned nonlinear suffix;
we do not replace that suffix with a linear approximation.

We run two tests with 12 independently trained models per target family:

\begin{itemize}
  \item \textbf{Does the certificate predict failure?} We fix the observer
  response at $g_E=0.85$ and vary the response mismatch $\rho$, initial bias,
  and target offset. Small positive and negative edits measure the actual target
  response $g_T$. We then ask whether Equation~\ref{eq:certificate} predicts
  which rollouts miss the target and by how much.

  \item \textbf{Does the interaction coordinate improve control?} We compare
  the natural first-order and pairwise estimator--direction pairs. We use
  $\gamma=1.15$ when the target contains the interaction and $\gamma=0$ as the
  additive control.
\end{itemize}

Integrated squared tracking error counts error across the full trajectory, not
only at the final step, and uncertainty is measured across the 12 independently
trained models, not across the many conditions within each model. None of the
main comparisons clips. Every natural comparison and 99.2\% of controlled
conditions move no farther than the largest distance among the four clean task
archetypes. The edits therefore have a familiar scale, although four archetypes
cannot prove that edited states remain on the full activation manifold.

The suffix is measurably nonlinear. At the largest offsets, its mean deviation
from the affine prediction is 16.1\% of output movement, and the certificate
residual is 20.2\% of mean absolute final error. We therefore ask whether the
affine quantities remain predictive, not whether the Transformer is linear.

They remain predictive. The main checks agree:

\begin{itemize}
  \item The raw certificate correlates $0.952$ with actual final error and gets
  the sign right in 96.0\% of conditions.
  \item Its residual MAE is 16.5\% of mean absolute actual error.
  \item After fitting only an intercept and slope on the other 11 seeds, its
  leave-one-seed-out MAE is $0.063$, compared with $0.302$ for the bias-only
  model and $0.367$ for the observer-pole-only model. It wins on every held-out seed.
\end{itemize}

Across this checked operating range, the three certificate terms capture most
of the error. This does not show that the suffix or its activation manifold is
locally linear.

\begin{table}[t]
\centering
\scriptsize
\begin{tabular}{@{}lrrrr@{}}
\toprule
Comparison & A & B & A--B [95\% CI] & Seeds \\
\midrule
Bias vs. full cert. & 0.302 & 0.063 & 0.238 [0.230, 0.246] & 12/12 \\
Pole vs. full cert. & 0.367 & 0.063 & 0.303 [0.294, 0.312] & 12/12 \\
FO vs. lifted, $\gamma=1.15$ & 2.632 & 1.870 & 0.762 [0.424, 1.101] & 11/12 \\
FO vs. lifted, $\gamma=0$ & 1.693 & 1.662 & 0.031 [-0.021, 0.103] & 6/12 \\
\bottomrule
\end{tabular}
\caption{Study-1 nonlinear-suffix results. A--B is positive when the full certificate or lifted pair has lower error. Intervals resample the 12 training seeds.}
\label{tab:phase5-nonlinear-suffix}
\end{table}


The direct control comparison gives the same message. The pairwise model
reduces integrated squared error by 29.0\% when the target contains the
``both features present'' term, with improvement in 11 of 12 independently
trained models. It gives no reliable advantage for the additive target. The
exact losses and intervals are reported in Table~\ref{tab:phase5-nonlinear-suffix}.

This result concerns the estimator and edit direction together. The mismatched
estimator--direction pairs are unstable in 10 of 12 interactional models and
also fail the action- and activation-scale checks. We therefore do not use them
to rank the estimators independently of direction. Appendix
Table~\ref{tab:nonlinear-suffix-factorial} reports all four cells, including
their self-response and convergence factor $p$.

\begin{figure}[H]
  \centering
  \includegraphics[width=0.98\linewidth]{nonlinear_suffix_certificate_and_control.png}
  \caption{Learned nonlinear-suffix experiment. The certificate combines
  starting mismatch, convergence factor, and response mismatch (left). The
  pairwise estimator--direction pair reduces control error for the
  interactional target but not for the additive target (right). Intervals
  resample 12 independently trained models.}
  \label{fig:nonlinear-suffix}
\end{figure}

\FloatBarrier
The study provides a bridge to pretrained models: a simple control-error
certificate can remain useful after nonlinear computation when edit scale and
loop stability are checked. It does not establish token-level or large-model
control. Closed-loop control evaluates an observer over a trajectory. What if
the decision instead selects one static intervention from a fixed candidate
set using predicted finite effects?
